%! Author = sriadityavedantam
%! Date = 3/25/26

\lesson{35}{Monday 3 November 2025 9:10}{Day 35}

\begin{theorem}{
    Suppose $f,g : (a,b] \to\rr$ are differentiable, $g(x)\to \infty$ as $x\to a^+$, and $g\uptick(x)\neq 0$ for all $x\in(a,b]$.
    If
    \[
        \lim_{x\to a^+}\dfrac{f\uptick(x)}{g\uptick(x)} = L,
    \]
    then
    \[
        \lim_{x\to a^+}\dfrac{f(x)}{g(x)} = L.
    \]
}
	Let $0 < \eps < 1$ be given.
    Choose $\delta > 0$ such that
    \[
        0 < x-a < \delta \implies \abs{\dfrac{f\uptick(x)}{g\uptick(x)} - L} < \eps.
    \]
    Fix $y\in(a,a+\delta)$.
    Then
    \[
        \abs{\dfrac{f\uptick(y)}{g\uptick(y)} - L} < \eps.
    \]
    For $a < x < y$, the Cauchy Mean Value Theorem applied on $[x,y]$ implies there exists $c = c_{xy}$ such that $x < c < y$ and
    \[
        f\uptick(c)(g(y) - g(x)) = g\uptick(c)(f(y) - f(x)).
    \]
    Hence
    \[
        \dfrac{f\uptick(c)}{g\uptick(c)} = \dfrac{f(y) - f(x)}{g(y) - g(x)} = \dfrac{f(x) - f(y)}{g(x) - g(y)}.
    \]
    Since $x<c<y<a+\delta$, we have
    \[
        \abs{\dfrac{f(x) - f(y)}{g(x) - g(y)} - L} < \eps.
    \]
    Define
    \[
        \alpha(x) \coloneqq \dfrac{f(x)}{g(x)},
        \qquad
        \beta(x) \coloneqq \dfrac{f(x)-f(y)}{g(x)-g(y)}.
    \]
    Since $g(x)\to\infty$ as $x\to a^+$, there exists $\delta\uptick > 0$ such that $0 < x-a < \delta\uptick$ implies
    \[
        \abs{\dfrac{f(y)}{g(x)-g(y)}} < \eps
        \qquad\text{and}\qquad
        \abs{\dfrac{g(y)}{g(x)-g(y)}} < \eps.
    \]
    For such $x$, we have
    \begin{align*}
        \abs{\alpha(x) - \beta(x)}
        &= \abs{\dfrac{f(x)}{g(x)} - \dfrac{f(x)-f(y)}{g(x)-g(y)}}\\
        &= \abs{\dfrac{f(y)}{g(x)-g(y)} - \dfrac{f(x)}{g(x)}\dfrac{g(y)}{g(x)-g(y)}}\\
        &\leq \abs{\dfrac{f(y)}{g(x)-g(y)}} + \abs{\dfrac{f(x)}{g(x)}}\abs{\dfrac{g(y)}{g(x)-g(y)}}\\
        &\leq \eps + \eps\abs{\alpha(x)}.
    \end{align*}
    Thus
    \begin{align*}
        \abs{\alpha(x)-L}
        &\leq \abs{\beta(x)-L} + \abs{\alpha(x)-\beta(x)}\\
        &\leq \eps + \eps + \eps\abs{\alpha(x)}\\
        &= 2\eps + \eps\abs{\alpha(x)}.
    \end{align*}
    Also
    \[
        \abs{\alpha(x)} \leq \abs{\alpha(x)-L} + \abs{L}.
    \]
    So
    \begin{align*}
        \abs{\alpha(x)-L}
        &\leq 2\eps + \eps(\abs{\alpha(x)-L} + \abs{L})\\
        &= (2+\abs{L})\eps + \eps\abs{\alpha(x)-L}.
    \end{align*}
    Hence
    \[
        (1-\eps)\abs{\alpha(x)-L} \leq (2+\abs{L})\eps.
    \]
    Therefore
    \[
        \abs{\alpha(x)-L} \leq \dfrac{(2+\abs{L})\eps}{1-\eps}.
    \]
    Since the right-hand side tends to $0$ as $\eps\to 0$, it follows that
    \[
        \lim_{x\to a^+}\dfrac{f(x)}{g(x)} = L.
    \]
\end{theorem}